% Independently audited complete proof; amendments make descent points explicit.
\subsubsection*{Statement B.2s (Proved): weak stabilization when a regular birational model exists}
\begin{theorem}
\label{B:B2-regular-model-stabilization}
Let $k$ be a field of characteristic $p>0$, and let $K/k$ be a
finitely generated regular extension.  Suppose that $K/k$ has a
proper regular model $Y$.  Given any proper geometrically normal
model $X_*$ of $K/k$, there is a proper geometrically normal
modification $X_0\to X_*$ such that every proper geometrically
normal modification $W\to X_0$ induces an isomorphism
\[
 \operatorname{Pic}^{\tau}_{X_0/k}
       \xrightarrow{\ \sim\ }\operatorname{Pic}^{\tau}_{W/k}.
\]
The assertion concerns an upper cone and does not assert that this
cone is cofinal in the entire model category.
\end{theorem}

\begin{proof}
Write $P=\operatorname{Pic}^{\tau}_{Y/k}$.  This is a commutative
group scheme of finite type; it need not be proper over $k$.
For every permitted model $X\to X_*$, normalize the main graph
between $X$ and $Y$, obtaining $T\to X,Y$.  Statement B.1g
identifies $\operatorname{Pic}^{\tau}_T$ with $P$.
Since the proper birational map $T\to X$ has universal direct
image of its structure sheaf equal to $\mathcal O_X$, pullback
gives a closed immersion
\[
                       P_X:=\operatorname{Pic}^{\tau}_{X/k}
                              \lhook\joinrel\longrightarrow P.
\]
Here $P_X$ is proper because $X$ is geometrically normal.
A further normal common graph shows that the immersion is
independent of the choice of $T$, and that these immersions are
compatible with domination.  We henceforth regard the $P_X$ as
closed subgroup schemes of $P$.

Choose $X_1\to X_*$ for which $\dim P_{X_1}$ is maximal.
Such a choice exists because these dimensions are nonnegative
integers bounded by $\dim P$.  Let $A_0$ be the largest abelian
subvariety of $P_{X_1}$, obtained from the stable multiplication
image in B.2r.  Since $P_{X_1}$ is proper, its geometric reduced
identity component is an abelian variety; multiplication by $p$
is surjective on that component.  The stable image therefore
has dimension $\dim P_{X_1}$, so $P_{X_1}/A_0$ is finite.
Every subsequent $X\to X_1$ has the same largest abelian
subvariety $A_0$ inside $P$: the inclusion of these abelian
subvarieties has equal dimensions, again using properness.
Statement B.2r shows that $P$ has no geometric torus.

We construct, after finitely many further permitted modifications,
a closed subgroup $H\subset P$ containing all the subsequent
$P_X$ and having maximal abelian subvariety exactly $A_0$.
Initially take $H=P$.  The largest abelian subvariety $B\subset H$
exists by the stable-multiplication-image construction in B.2r.
It contains $A_0$.  If it equals $A_0$, the construction is finished.
Suppose instead that $\dim B>\dim A_0$.

\paragraph{A bad divisor attached to the excessive abelian part.}
The Albanese torsor of the proper geometrically integral variety
$Y$ is dual to the maximal abelian subvariety of its Picard
scheme \cite[Theorem 3.2 and Remark 3.3, pp.~12--13 of the
author manuscript]{C_ACMV2026}.  Pushing this torsor out along
the dual of the inclusion of $B$ gives a torsor $\mathcal A_B$
under $B^\vee$ and a morphism $Y\to\mathcal A_B$ whose pullback
on $\operatorname{Pic}^0$ is the inclusion $B\subset P$.
Thus there is a rational map $X_1\dashrightarrow\mathcal A_B$.
Using the torsor avoids any extension of the ground field to
choose an origin.

Let $G$ be the normalization of its closed graph.  This is a
proper normal model with maps to $X_1$ and $\mathcal A_B$.
If $G$ were geometrically normal, it would be a permitted
modification of $X_1$ and its Picard image would contain $B$,
contrary to the maximality of $\dim A_0$.
Consequently $G$ is not geometrically normal.
Normality implies geometric $S_2$ under field extensions
\cite[Tags 045M and 0339]{StacksB}.  Geometric integrality
follows from the regularity of $K/k$.  By Serre's criterion,
there is a codimension-one point of $G$ whose local ring fails
geometric normality.  More explicitly, the failure occurs after
a finite algebraic extension, and a bad codimension-one point
there lies above a codimension-one point of $G$.
Denote the resulting divisorial valuation of $K$ by $v$ and its
residue field by $\kappa(v)$.

This valuation cannot have a divisorial centre on any permitted
model.  Otherwise the common normal graph over that model and
$G$ would be an isomorphism near the two codimension-one DVRs,
identifying the bad local ring with a geometrically normal one.
Its centre therefore has dimension at most $n-2$, where
$n=\operatorname{trdeg}_k K$.
The divisor on $G$, on the other hand, has image of dimension
$n-1$ in $X_1\times\mathcal A_B$, because normalization of
the graph is finite over its image.

Among permitted modifications of $X_1$, choose $X_2$ on which
the centre $C_2$ of $v$ has maximal dimension.  For every
permitted $W\to X_2$, its centre $C_W$ maps dominantly and
properly to $C_2$ and has the same dimension.  In particular
\[
                 F:=k(C_2)\ \subset\ k(C_W)
\]
is a finite extension inside the fixed residue field $\kappa(v)$.
We use the reduced centre in this assertion.

\paragraph{Restriction to a fixed proper generic fibre.}
Take a normal proper common graph $T$ dominating $G,Y,X_2$.
Let $D$ be the normalization of the strict transform on $T$
of the divisor defining $v$.  It is proper over the centre
$C_2$, and the generic fibre $D_F$ is a proper normal integral
$F$-scheme with function field $\kappa(v)$.
Pullback from $Y$ followed by restriction defines a homomorphism
\[
                P_F\longrightarrow\operatorname{Pic}_{D_F/F}.
\]
Let $K_F$ be its closed scheme-theoretic kernel.

There is a largest closed $k$-subgroup $K_v\subset P$ whose
base change to $F$ factors through $K_F$, with the corresponding
factorization property on every test scheme.
To see this directly, on an affine open $\operatorname{Spec}R$
of $P$ choose finitely many generators of the ideal of $K_F$
in $R\otimes_k F$.  Expand the finitely many coefficients in
a $k$-basis of $F$ and let $I\subset R$ be generated by their
coefficients in $R$.  For every $k$-algebra $S$, a map $R\to S$
annihilates $I$ exactly when its base change to $F$ annihilates
the original ideal.  Thus these closed subschemes glue; the
group law, inverse and identity restrict by this universal
property.  No Weil restriction along an infinite field
extension is being asserted.

Every $P_W$ for a permitted $W\to X_2$ is contained in $K_v$.
Indeed, take a further normal common graph over $T$ and $W$,
and normalize the strict transform $D'$ of $D$.  The morphism
$D'_F\to D_F$ is proper birational between normal integral
schemes.  Its direct image of the structure sheaf
is $\mathcal O_{D_F}$; this remains so on every parameter base
change, since every $F$-scheme is flat over the field $F$.
Thus pullback of relative Picard sheaves
is injective.  The map $D'_F\to W$ factors through
$\operatorname{Spec}k(C_W)$ at the generic fibre of $C_2$.
A line bundle pulled back from $W$ is therefore, on this
generic fibre, pulled back from this finite field extension of
$F$.  Its relative Picard class is zero on all parameter
schemes: a line bundle on a finite locally free algebra becomes
trivial after the smooth surjective generator-open cover used
in B.2m.  This argument includes nonreduced parameters.
Injectivity of pullback to $D'_F$ proves the assertion on $D_F$,
and the defining property of $K_v$ proves $P_W\subset K_v$.

The subgroup $B$ is not contained in $K_v$.
Indeed the image of $D$ in $X_1\times\mathcal A_B$ still has
dimension $n-1$, whereas $\dim C_2\le n-2$.
The image of $D$ in $C_2\times\mathcal A_B$ dominates this
first image and therefore has dimension at least $n-1$.
Consequently the restriction $D_F\to(\mathcal A_B)_F$ has
positive-dimensional image.  After an algebraic closure of
$F$, at least one reduced irreducible component has a
nonconstant map to the torsor, and so does its normalization.
It is a proper geometrically integral variety.  The Albanese
universal property and Picard duality show that a nonconstant
map from it to an abelian torsor has nonzero pullback on
$\operatorname{Pic}^0$.  Since this pullback is the restriction
of $B\to P_F$, we obtain $B\not\subset K_v$.

Replace $H$ by $H\cap K_v$ and the current model by $X_2$.
At the next step we relabel this current model as $X_1$;
the subgroup $A_0$ stays fixed throughout the construction.
All later Picard images lie in the new $H$.  Its largest
abelian subvariety is contained in the proper closed subgroup
$B\cap K_v$ of $B$, and hence has strictly smaller dimension.
It still contains $A_0$.  Repeating this construction can take
at most $\dim B-\dim A_0$ steps.  We therefore reach a permitted
model $X_3$ and a fixed closed subgroup $H\subset P$ such that
all $P_W$, for $W\to X_3$, lie in $H$ and its largest abelian
subvariety is $A_0$.

\paragraph{The remaining finite quotient has a uniform order bound.}
By B.2r there is an integer $N>0$ with $[N]H\subset A_0$.
For every later $W$, the quotient $P_W/A_0$ is finite, and
the map $P_W[N]\to P_W/A_0$ is fppf-surjective.
For an explicit verification, if $x$ is a local section of
$P_W$, then $Nx\in A_0$; fppf locally we can choose
$a\in A_0$ with $Na=Nx$, since multiplication by $N$ on an
abelian variety is faithfully flat.  The section $x-a$ belongs
to $P_W[N]$ and represents the same quotient class.

After one algebraically closed extension, every such $P_W$
embeds compatibly into a fixed $g$-dimensional Jacobian $J$ by
B.2d.  To apply that statement, one may use normal graphs over
one fixed projective model after algebraic closure; no
geometrically normal common refinement over $k$ is needed.
We obtain
\[
       |P_W/A_0|\le |P_W[N]|\le |J[N]|=N^{2g}.
\]
Choose $X_0\to X_3$ for which the finite order $|P_{X_0}/A_0|$
is maximal.  Every later closed immersion
$P_{X_0}\hookrightarrow P_W$ induces a closed immersion of
finite quotients of the same order, hence an isomorphism.
It follows that $P_{X_0}\to P_W$ itself is an isomorphism.
This proves the theorem without a filteredness assumption.
\end{proof}

\begin{corollary}
The weak upper-cone assertion holds over every field of positive
characteristic for function fields of transcendence degree at most
three, whenever the class of proper geometrically normal models
is nonempty.
\end{corollary}
\begin{proof}
For a geometrically integral function field in dimension at most
three, a proper regular birational model exists by
\cite[Theorem 1.1, pp.~3--4 of the author manuscript]{CossartPiltant2019B}.
The hypotheses of that theorem include finite-type schemes over
arbitrary fields: they are excellent, and the asserted resolution
is proper and birational.
If $K$ has nontrivial algebraic constants $k_0/k$, geometric
normality of a permitted model implies that $k_0/k$ is finite
separable.  We apply the result over $k_0$, where $K/k_0$ is
regular, and descend the resulting Picard isomorphisms through
$\operatorname{Res}_{k_0/k}$, as in B.2o.
\end{proof}
